% New writing for the opening of Volume II. This file is never overwritten by split_volumes.py.

This volume computes the objects that Volume I defines, and its centre is the Cauchy deconvolution approach to the Riemann Hypothesis. The conventions are those of Volume I (Section~\ref{ov1:sec:conventions}): every statement is proved, computed, certified, conditional, heuristic or open, and the hypothesis of a conditional statement is named. This volume writes the level $\theta$ and the density $u_\theta(v)=\pi^{-1}\re(\xi'/\xi)(\theta+iv)$, which Volume I writes as $v_\alpha/\pi$ with $\alpha=\theta$.

\section{The path to RH, seen from the computation}\label{ov2:sec:path}

Volume I (Section~\ref{ov1:sec:path}) argues that the clearest path to RH is a descent in level: the Thorin measure at level $\theta$ is positive exactly when $\zeta$ has no zero with $\re s>\theta$ (Theorem~\ref{cx:thm:positivity}), its global minimum sits at the origin at level one with no hypothesis (Proposition~\ref{bp:prop:v1min} of Volume I), and RH is the statement that it stays there at every level down to $\frac12$ (Theorem~\ref{thm:dbeta:minall} of Volume I). In this volume the same statement is a deconvolution. Under RH the level one density is the Cauchy smoothing of the zeros, $u_1=P_{1/2}*\mathfrak U_{1/2}$, and in general a positive deconvolution of $u_1$ down to level $\frac12+\varepsilon$ exists exactly when no zero has real part above $\frac12+\varepsilon$ (Theorem~\ref{tf:thm:deconv}). The level one density is computable from the Euler product; RH is the statement that the Cauchy flow can be run backwards from it to the critical line through positive measures.

Part II of this volume does that computation as well as it can be done, and the results show both its power and its limit.
\begin{enumerate}
\item The operator (Chapter~\ref{chap:cd}). Cauchy deconvolution is Hadamard's ill posed problem: unique when solvable, H\"older stable between levels and only logarithmically stable at the boundary, with a critical width beyond which no positive source exists. For the zeta family that width is one minus the supremum of the real parts of the zeros.
\item The backward flow (Chapters~\ref{chap:lo} and \ref{chap:cx}). From samples of the level one density alone, the flow run to level $0.75$ places a peak within $0.035$ of each of the $109$ zeros tested between heights $100$ and $300$; the EM algorithm of Shepp and Vardi does it without ringing, and a posteriori bounds certify window masses in ball arithmetic up to height $1419.9$.
\item The bottom of the flow (Chapter~\ref{chap:es}). At level $\frac12$ the deconvolution is Weil's explicit formula. For a general measure it is exponentially ill posed, heuristically requiring the primes up to $e^{\pi/\delta}$ for resolution $\delta$; for a finite sum of atoms it is Prony's problem and only polynomially ill posed, and the matrix pencil recovers ordinates and real parts from level one data alone. But it reads the real parts; it does not constrain them.
\item The descent itself (Chapters~\ref{chap:cf} and \ref{chap:bsy}). Particle filters run the flow in level, nested sampling scans heights, and the criterion of Balazard, Saias and Yor is solved by Monte Carlo at every level. Each works in a window of heights.
\end{enumerate}
The limit is the one Volume I names. RH is a statement for all heights, and a perturbation of $\xi$ that agrees with it below any computed height passes the finite tests of the heat trace listed there (Theorem~\ref{thm:ch13:perturb}, Chapter~\ref{vol2:ch1}). So the computations of this volume can refute RH by certifying a zero off the line, can certify that windows are clean, and can measure how ill posed each step of the descent is; they cannot take the step. What would close the gap is an analytic positivity argument that uses the Euler product at level one and carries it to the line, and the last chapter of this volume (Chapter~\ref{vol2:ch24}) collects what such an argument would need. Statements that assume the zero free half plane $\re s>\frac78$ announced in October 2026, which has not been refereed, are marked where they occur, in particular in Chapters~\ref{vol2:ch22} and~\ref{vol2:ch23}.

\section{What a finite computation can decide}\label{ov2:sec:finite}

The rest of this chapter makes the limit precise: what a computation of the Thorin measure can reach, what a finite set of positivity tests cannot decide, and how computations of the zeros reached their present height.
